\subsection{图片检测模块}

\paragraph*{业务流程}

\paragraph*{用例图}

\subsubsection{功能点——提取图片（FR-TPJC-0001）}

\paragraph*{简述}

从用户上传的 PDF、压缩包及论文文件等非图像格式文件中自动抽取学术图像，便于后续进行图像AI鉴伪检测。

该功能旨在自动化处理多种文件格式，确保图像检测流程顺利进行。

系统能够快速提取高质量图像数据，为后续AI算法提供输入。

\paragraph*{前提条件}

用户上传文件格式为 PDF、ZIP、RAR、DOCX、TAR、7Z 等非图像格式。

文件大小不超过系统规定上限（如 500MB）。

文件中至少包含一张图像。

\paragraph*{主要流程}

\begin{enumerate}
    \item 用户上传文件；
    \item 系统自动识别文件格式；
    \item 调用对应解析工具：
    \begin{enumerate}
        \item PDF解析库
        \item 压缩文件解压工具
        \item DOCX文档解析工具
    \end{enumerate}
    \item 对PDF逐页提取嵌入图像；
    \item 对压缩文件递归遍历子目录提取图像；
    \item 对论文文档抽取正文插图、实验图及附录图像；
    \item 图像预处理：
    \begin{enumerate}
        \item 格式统一为 JPG / PNG
        \item 分辨率标准化
        \item 色彩空间统一
    \end{enumerate}
    \item 临时存储；
    \item 生成图像文件列表。
\end{enumerate}

\paragraph*{后继结果}

系统输出结构化图像列表，包含：

\begin{enumerate}
    \item 文件名
    \item 格式
    \item 分辨率
    \item 来源页码
    \item 图像编号
\end{enumerate}

图像列表以 JSON 格式存储。

\paragraph*{补充说明}

支持批量文件处理；

支持异常提示：

文件损坏、格式不支持、解析失败。

\subsubsection{功能点——学术图像AI鉴伪（FR-TPJC-0002）}

\paragraph*{简述}

基于深度学习和图像处理技术，对提取图像进行自动化学术造假检测。

检测内容包括：

\begin{enumerate}
    \item 美化修饰
    \item 背景掩盖
    \item 复制伪造
    \item 跨图拼接
    \item AI生成图像识别
\end{enumerate}

系统能够高效识别图像异常区域并输出可解释结果。

\paragraph*{前提条件}

系统已成功提取图像数据；

图像已完成预处理。

\paragraph*{主要流程}

\begin{enumerate}
    \item 图像预处理：
    \begin{enumerate}
        \item 尺寸归一化
        \item 色彩归一化
        \item 噪声去除
    \end{enumerate}
    
    \item 特征提取：
    \begin{enumerate}
        \item CNN深度特征
        \item 边缘纹理特征
        \item 频域特征
    \end{enumerate}
    
    \item 图像鉴伪检测：
    \begin{enumerate}
        \item 美化修饰检测
        \item 背景掩盖检测
        \item 同图复制检测
        \item 跨图拼接检测
        \item AI生成痕迹检测
    \end{enumerate}
    
    \item 计算风险评分；
    
    \item 输出可疑区域坐标；
    
    \item 生成检测报告。
\end{enumerate}

\paragraph*{后继结果}

系统输出检测报告，包括：

\begin{enumerate}
    \item 原图
    \item 标注图
    \item 可疑区域说明
    \item 造假类型
    \item 风险等级
    \item 置信度评分
\end{enumerate}

报告支持 JSON / PDF / HTML 格式。

\paragraph*{补充说明}

支持单图检测与批量检测；

支持综合鉴伪报告生成；

支持可视化展示。

\subsubsection{功能点——图像可疑区域细粒度标注（FR-TPJC-0003）}

\paragraph*{简述}

对AI检测出的疑似造假区域进行细粒度可视化标注，辅助用户和人工审核人员快速定位异常区域。

\paragraph*{前提条件}

学术图像AI鉴伪任务已完成。

\paragraph*{主要流程}

\begin{enumerate}
    \item 获取AI检测输出的可疑区域坐标；
    \item 在原图上绘制高亮边框；
    \item 对不同造假类型使用不同标注标签；
    \item 支持局部放大查看；
    \item 展示区域风险评分。
\end{enumerate}

\paragraph*{后继结果}

输出标注图像，包括：

\begin{enumerate}
    \item 异常区域框选
    \item 区域编号
    \item 风险等级
    \item 类型说明
\end{enumerate}

支持用户下载标注图。

\paragraph*{补充说明}

支持多区域同时标注；

支持细粒度像素级区域高亮；

支持人工审核二次标记。

\subsection{学术内容检测模块}

\paragraph*{迭代说明（职责边界）}

本节中标识为 FR-LWJC-xxxx 与 FR-NRJC-0002 的需求，在**用户端（前端）**侧重检测结果的**展示、导航与可视化交互**：包括列表与详情布局、段落/图表联动、概率与结论文本的呈现、加载与错误状态等。全篇 AIGC 判定、段落级识别与概率计算、事实性推理及解释文案所依赖的**模型与算法在后端执行**；前端根据接口返回的结构化数据与任务状态进行渲染，不在此需求中规定具体模型与实现细节。

\subsubsection{功能点——全篇论文AIGC检测（FR-LWJC-0001）}

\paragraph*{简述}

在用户完成论文类检测任务后，前端展示“全篇层面”的 AIGC 检测结论与摘要指标（如全文风险等级、AI 贡献占比或等价字段，以后端约定为准），使用户快速了解整篇论文的自动鉴伪概况。

\paragraph*{前提条件}

用户已登录；对应论文检测任务已创建且后端已返回可供展示的全文级结果（或返回处理中/失败状态码）。

\paragraph*{主要流程}

\begin{enumerate}
    \item 用户在检测任务详情或论文检测专页进入“AIGC 全文检测”结果区。
    \item 前端根据任务 ID 请求全文级结果接口，展示摘要卡片（结论、关键数值、更新时间等）。
    \item 若任务仍在计算中，展示进度或排队状态，并在就绪后自动或手动刷新。
    \item 用户可从摘要区跳转到段落级明细（与 FR-LWJC-0002、0003 联动，若接口已提供锚点或段落 ID）。
\end{enumerate}

\paragraph*{后继结果}

用户获得可理解的全文级结论；接口异常或超时时展示错误提示与重试入口。

\subsubsection{功能点——AI生成段落识别（FR-LWJC-0002）}

\paragraph*{简述}

前端将后端标识的疑似 AI 生成段落以列表或文中高亮等形式展示，支持用户按段落查看系统标注的位置与简要标签，便于对照原文审阅。

\paragraph*{前提条件}

论文检测任务已产生段落级标注数据；用户具备该任务的查看权限。

\paragraph*{主要流程}

\begin{enumerate}
    \item 前端拉取段落级标注列表（含段落序号、起止位置或文本片段引用等）。
    \item 在正文视图或侧栏列表中高亮对应段落；用户点击某条标注时，视图滚动至该段落。
    \item 可选展示每条标注的简短类型或置信度区间（以后端字段为准）。
\end{enumerate}

\paragraph*{后继结果}

用户可定位全部被标注段落；无标注时展示说明“未检出疑似 AI 段落”或等价文案。

\subsubsection{功能点——段落级AI生成概率分析（FR-LWJC-0003）}

\paragraph*{简述}

针对被识别或全量段落，前端展示后端给出的 AI 生成概率或分档结果（如数值、等级、热力示意），帮助用户理解模型在段落粒度上的判断强度。

\paragraph*{前提条件}

后端对相应段落返回概率或分档字段；前端具备可视化组件（表格、条形、颜色映射等）展示能力。

\paragraph*{主要流程}

\begin{enumerate}
    \item 用户在段落视图或独立“概率分析”面板中查看各段概率或分档。
    \item 支持按概率排序、筛选高于某阈值的段落（若产品启用）。
    \item 与 FR-LWJC-0002 联动：点击段落时同步显示该段概率。
\end{enumerate}

\paragraph*{后继结果}

概率数据与段落一一对应展示；缺失字段时对应单元显示“—”或隐藏该项并附说明。

\subsubsection{功能点——事实性鉴伪子结论生成（FR-LWJC-0004）}

\paragraph*{简述}

前端展示后端生成的事实性鉴伪**子结论**列表或结构化条目（如“数据矛盾”“引用不实”等类别及条目化描述），使用户了解除 AIGC 以外的事实核查维度上的发现。

\paragraph*{前提条件}

该检测任务启用了事实性鉴伪能力且后端已返回子结论结构；用户有查看权限。

\paragraph*{主要流程}

\begin{enumerate}
    \item 前端请求事实性子结论接口，按类别或严重级别分组展示。
    \item 用户展开单条子结论查看摘要；若后端提供关联原文位置或引用，前端提供跳转或悬浮展示。
\end{enumerate}

\paragraph*{后继结果}

子结论完整可读；若无子结论则展示空状态或“未触发事实性规则”类说明。

\subsubsection{功能点——事实性鉴伪原因说明（FR-LWJC-0005）}

\paragraph*{简述}

对每条事实性鉴伪子结论，前端展示后端附带的**原因说明**与依据要点（如触发的规则说明、证据片段引用等），以提升可理解性；文案内容由后端生成，前端负责排版与折叠展开。

\paragraph*{前提条件}

子结论记录中包含原因说明字段或关联证据结构。

\paragraph*{主要流程}

\begin{enumerate}
    \item 用户在子结论卡片中展开“原因说明”区域。
    \item 若存在多条证据，按列表或步骤展示；支持复制文本（若产品允许）。
\end{enumerate}

\paragraph*{后继结果}

用户能理解单条结论的判定逻辑表层信息；超长文本支持折叠与“展开全文”。

\subsubsection{功能点——同行评审Review检测（FR-PLJC-0001）}

\subsubsection{功能点——模板化评审倾向识别（FR-PLJC-0002）}

\subsubsection{功能点——综合鉴伪结果生成（FR-NRJC-0001）}

\subsubsection{功能点——检测结果可视化解释（FR-NRJC-0002）}
